% ============================================================================
%  Cover page - Technical Design Document
% ----------------------------------------------------------------------------
%  Requires (already loaded in main.tex): graphicx, xcolor, geometry.
%  Edit the \cover... fields below to update title, subtitle, version and the
%  team list; the layout reflows automatically.
%
%  Typography: the page stays in the document's roman (serif) face rather than
%  switching to sans, so the cover reads as part of a formal engineering report
%  instead of a slide. Hierarchy comes from size and letterspacing, not weight
%  alone.
%
%  Logos live in images/ alongside the other document graphics:
%      images/SPACECHALLENGES.png  (program)
%      images/armonia_logo.png     (team)
%  Both source files are square (4500x4500 and 1024x1024), so they are drawn
%  into identical square boxes of side \coverlogosize and therefore match
%  exactly. To resize both, change \coverlogosize alone. To trim one logo's
%  own internal whitespace, give it the optional scale argument, e.g.
%  \coverlogo[1.1]{...} - the box stays the same size, only the artwork inside
%  it grows. Each logo is guarded: a missing file leaves blank space of the
%  same size rather than failing the build.
% ============================================================================

\definecolor{coverrule}{HTML}{2E75B6}  % divider rules (matches ctcol accent)
\definecolor{covermeta}{HTML}{404040}  % metadata label grey

% ---------- Editable cover fields ----------
\newcommand{\coverTitle}{Technical Design Document}
\newcommand{\coverSubtitle}{Cubli}
\newcommand{\coverTagline}{A Three-Axis Reaction-Wheel Balancing Cube}
\newcommand{\coverVersion}{v1.0}
\newcommand{\coverOrganization}{Space Challenges Program 2026}

% Team list, laid out as an explicit three-column grid: names separated by & and
% rows by \\. Keeping it explicit means the columns line up exactly; to add a
% member, extend a row or start a new one.
\newcommand{\coverTeam}{%
    Athanasia Nikolova & Andrea Liang & Deyan Nikolaev Vlaev \\[0.85em]
    Niccolo & Suvanna Viriyanti Wu & Pablo Urioste \\}

% ---------- Helpers ----------
% Full-width horizontal divider; argument = rule thickness.
\newcommand{\coverrule}[1]{{\color{coverrule}\rule{\linewidth}{#1}}}

% Letterspaced small-caps run, used for the label above the title and for the
% metadata keys. Gives the cover a formal, engraved feel without extra packages.
\newcommand{\coverspaced}[1]{%
    {\scshape\fontsize{10}{12}\selectfont
     \makebox[\linewidth][c]{\color{covermeta}#1}}}

% ============================ LOGO TUNING ==================================
% The two logo files carry different amounts of built-in white space around the
% artwork, so drawing them at the same nominal size does NOT make them look the
% same size or sit on the same line. The six knobs below exist to correct that
% by eye. Each logo has its own scale and its own horizontal/vertical offset,
% and they are completely independent - changing one never moves the other.
%
%   ...width  : how wide the artwork is drawn. Bigger value = bigger logo.
%               Use this to make the two marks look optically equal.
%   ...raise  : vertical nudge. POSITIVE moves the logo UP, negative moves DOWN.
%   ...shift  : horizontal nudge. POSITIVE moves the logo RIGHT, negative LEFT.
%
% Tune in small steps (0.05cm-0.1cm); recompile and look. Nothing else on the
% cover depends on these values.
% ---------------------------------------------------------------------------

% ----- Left logo: Space Challenges -----
\newlength{\sclogowidth}   \setlength{\sclogowidth}{3.30cm}
\newlength{\sclogoraise}   \setlength{\sclogoraise}{0cm}
\newlength{\sclogoshift}   \setlength{\sclogoshift}{0cm}

% ----- Right logo: Armonia -----
\newlength{\armlogowidth}  \setlength{\armlogowidth}{3.05cm}
\newlength{\armlogoraise}  \setlength{\armlogoraise}{0cm}
\newlength{\armlogoshift}  \setlength{\armlogoshift}{-1cm}

% Gap between the two logos.
\newlength{\coverlogogap}  \setlength{\coverlogogap}{2.2cm}

% Height reserved for the logo row. It is fixed and independent of the logos,
% so raising or lowering either mark never shifts the title block underneath.
% Increase it if you enlarge a logo enough to overflow the row.
\newlength{\coverlogorow}  \setlength{\coverlogorow}{3.4cm}

% \coverlogo{<width>}{<raise>}{<shift>}{<path>} - place one logo at the given
% width, nudged by raise (up) and shift (right). The box carries no height of
% its own, so a vertical nudge moves the artwork without disturbing the height
% of the logo row or anything below it.
\newcommand{\coverlogo}[4]{%
    \IfFileExists{#4}
        {\hspace*{#3}\raisebox{#2}[20pt][-20pt]{%
            \includegraphics[width=#1]{#4}}}
        {}}
% ===========================================================================

% Aligned key-value row of the metadata block.
\newcommand{\covermetarow}[2]{%
    \hfill\coverspacedlabel{#1} & #2 \\[0.5em]}
\newcommand{\coverspacedlabel}[1]{{\scshape\color{covermeta}#1}}

\begin{titlepage}
    \centering
    % titlepage appends its own \vfill, which would push the trailing metadata
    % and rule onto a discarded overflow page. Zeroing \parskip and holding the
    % whole cover in one top-anchored block keeps every element on this sheet.
    \setlength{\parskip}{0pt}

    % ---------- Logo row: program logo left, team logo right ----------
    % Each logo is positioned by its own width/raise/shift knobs defined in the
    % LOGO TUNING block above. \coverlogorow reserves a fixed height for the
    % row, so nudging a logo up or down never moves the title block below it.
    \vspace*{0.2cm}
    \makebox[\linewidth][c]{%
        \rule{0pt}{\coverlogorow}%
        \coverlogo{\sclogowidth}{\sclogoraise}{\sclogoshift}%
                  {images/SPACECHALLENGES.png}%
        \hspace{\coverlogogap}%
        \coverlogo{\armlogowidth}{\armlogoraise}{\armlogoshift}%
                  {images/armonia_logo.png}%
    }

    \vspace{1.5cm}

    % ---------- Title block ----------
    \coverspaced{Space Challenges Program 2026}

    \vspace{0.4cm}
    \coverrule{1.1pt}
    \vspace{0.9cm}

    {\fontsize{28}{34}\selectfont\scshape \coverTitle\par}

    \vspace{0.9cm}
    \coverrule{0.8pt}

    \vspace{1.2cm}

    % ---------- Project identity ----------
    {\fontsize{44}{50}\selectfont\bfseries \coverSubtitle\par}
    \vspace{0.55cm}
    {\large\itshape \coverTagline\par}

    % Explicit gaps rather than \vfill: inside titlepage the surrounding glue
    % absorbs the stretch, so fixed spacing is what actually holds the project
    % name high on the page and keeps the lower block from bunching.
    \vspace{2.4cm}

    % ---------- Team block ----------
    % Names sit on a three-column grid under a small-caps heading, so the list
    % reads as a deliberate block rather than a run of comma-separated text.
    \coverspaced{Project Team}
    \vspace{0.6cm}

    \begin{minipage}{0.92\textwidth}
        \centering
        \large
        \begin{tabular}{@{}*{3}{>{\centering\arraybackslash}p{0.32\linewidth}}@{}}
            \coverTeam
        \end{tabular}
    \end{minipage}

    \vspace{1.5cm}

    % ---------- Document metadata ----------
    \begin{minipage}{0.62\textwidth}
        \renewcommand{\arraystretch}{1.0}
        \begin{tabular}{@{}p{0.42\linewidth}@{\hspace{1.4em}}l@{}}
            \covermetarow{Organization}{\coverOrganization}
            \covermetarow{Document version}{\coverVersion}
            \covermetarow{Date of issue}{\today}
        \end{tabular}
    \end{minipage}

    \vspace{0.9cm}
    \coverrule{1.1pt}
    \vspace*{0.2cm}

\end{titlepage}
